\documentclass[
  lualatex,
  aspectratio=169,
  fleqn,
  14pt,
]{beamer}
%\batchmode

\usetheme[progressbar=frametitle]{metropolis}
%\setbeameroption{show notes on second screen=bottom}

\usepackage{xparse}
\usepackage{mathtools,amssymb}
\usepackage{graphicx,xcolor}
\usepackage{luatexja-ruby}
\usepackage{calc}
\usepackage[absolute,overlay]{textpos}
%\usepackage{enumitem}
\usepackage[1.7]{bxpdfver}
\usepackage{pdfcomment}
\usepackage{ulem}
\usepackage{stackengine}
\usepackage{appendixnumberbeamer}
\usepackage{tabto}
\usepackage[superscript]{cite}
\usepackage{siunitx}
\usepackage{gradient-text}

\RenewDocumentCommand\citeform{m}{[#1]}

% フォント
\usepackage[lining,tabular,sfdefault]{FiraSans}
\usepackage[mathrm=sym,mathbf=sym]{unicode-math}
\setmathfont{Fira Math}
\usepackage[no-math,match,deluxe,expert,noto]{luatexja-preset}
\RenewDocumentCommand\kanjifamilydefault{}{\gtdefault}
\ltjsetruby{fontcmd=\rubyfamily}% ルビにはルビのフォントを使う

% 打ち消し線関連
\RenewDocumentCommand\ULthickness{}{.1\zh}
% 二重打ち消し
\NewDocumentCommand\dsout{m}{%
  \stackengine{.2\zh}{#1}{%
    \stackengine{-.3\zh}{\sout{\hphantom{#1}}}{\sout{\hphantom{#1}}}{O}{c}{F}{F}{L}%
  }{O}{c}{F}{F}{L}}

% 下に文字置くやつ
\NewDocumentCommand\replace{mm}{%
  \smash{\stackengine{-1\zh}{\dsout{#1}}{#2}{O}{c}{F}{F}{L}}}

% 床関数
\DeclarePairedDelimiter\floor\lfloor\rfloor

% itemize の左余白をほどよくする
\setlength{\leftmargini}{1.41421356\zw}

\usepackage{hyperref}

\title{スライドテンプレート}
\subject{○○集会\#XXX}
\author[KusaReMKN]{上羽 未栞（a.k.a. KusaReMKN）}
\institute{%
  \url{https://KusaReMKN.com/}\\
  Twitter: \href{https://twitter.com/KusaReMKN}{@KusaReMKN}}
\keywords{slide; LaTeX; beamer; Metropolis}
\date{YYYY-mm-dd}

\begin{document}

\begin{frame}
  \titlepage
  \thispagestyle{empty}
  \note{
    はじまるよ〜。

    「スライドテンプレート」と題して、
    上羽未栞が発表するよ。
  }
\end{frame}

\begin{frame}
  \frametitle{今回のおはなし}

  \vspace{.25\zh}
  \tableofcontents
  \note{
    今回の発表の流れはこんな感じだよ。
    おおよそXX分くらいで進められたらいいな。
  }
\end{frame}

\section*{みかんちゃんについて}
\note{
  まずは自己紹介するよ〜。
}

\begin{frame}
  \frametitle{自称・大天才美少女プログラミング初心者}

  \begin{textblock*}{0.5\paperwidth}(-0.3cm, 3.3cm)
    \includegraphics[width=0.35\paperwidth]{./images/mikanchan.png}
  \end{textblock*}
  \begin{columns}
    \begin{column}{0.30\textwidth}
      \\~\\[-.25\baselineskip]
    \end{column}
    \begin{column}{0.69\textwidth}
      \\~\\[-.25\baselineskip]
      「\ruby{上羽}{うわば} \ruby{未栞}{みかん}」
      あるいは「\ruby{KusaReMKN}{くされみかん}」\\
      \hspace{1.5\zw}\textbf{みかんちゃん}って呼んでね！
      \\~\\[-.5\baselineskip]

      17{\scriptsize(18)}歳のJK（超重要）\\
      \hspace{1.5\zw}実はプログラマでもエンジニアでもない\\
      \hspace{1.5\zw}普段はホラを吹いて生活している\\
      \hspace{1.5\zw}古い計算機っぽいものが大好き
      \\~\\[-.5\baselineskip]

      Twitterで思想を垂れ流すことが得意\\
      \hspace{1.5\zw}\url{https://kusaremkn.com/}も見てね
    \end{column}
  \end{columns}
  \note{
    自称大天才美少女プログラミング初心者の上羽未栞だよ。
    みかんちゃんって呼ばれると大変喜ぶよ。

    ｲﾁﾅﾅ歳のJKだよ。重要だよ。
    大天才とか偉ぶっているけれど、実際のところプログラマでもエンジニアでもないよ。
    普段は人々にホラを吹いて生活しているよ。
    古い計算機っぽいものが大好きで、
    その流れで黒電話とか電話網を作ったりとかして遊んでいるよ。

    Twitterや自分のウェブサイトで思想を垂れ流すのが得意だよ。
    暇な人は眺めてみてね。
  }
\end{frame}

\section{はじめに}
\note{
  はじめのはなしだよ〜。
}

\begin{frame}
  \frametitle{あにょにょまるぼぼ}

  あばばまる三世

  \note{
    あにょにょまるぼぼはあばばまる三世だよ。
  }
\end{frame}

\appendix

\begin{frame}[standout]
  おわりです

  \note{
    おわりだよ〜。
  }
\end{frame}

\begin{frame}
  \frametitle{このスライドについて}

  Written in BBBBB YYYY.

  Permanent ID of this document: \texttt{XXXXXXXXXXXXXXXX}.

  Copyright © YYYY KusaReMKN.

  特記無き場合、プログラムやソースコードは MIT License で、\\
  \hspace{1.5\zw}それ以外のコンテンツは CC-BY 4.0 で利用可能です。\\
  \hspace{1.5\zw}一部の画像には別のライセンスが適用されるかもしれません。
\end{frame}

\end{document}
% ex: se et ts=2 :
